\documentclass{wlscirep}

\newcommand{\fcgr}{Fc\textgamma{}R}

\addbibresource{Library.bib}

\title{Dissecting \fcgr{} Regulation Through a Multivalent Binding Model}

\author[1]{Ryan A. Robinett}
\author[1,*]{Aaron S. Meyer}
\affil[1]{Koch Institute for Integrative Cancer Research, Massachusetts Institute of Technology, Cambridge, MA 02139}

\affil[*]{a@ameyer.me}

\keywords{Immune complex, multivalent ligand, receptor cross-linking, avidity, immunoglobulin G, Fc\textgamma{} receptors}

\begin{abstract}
Many immune receptors transduce activation across the plasma membrane through their clustering. In the case of Fc\textgamma{} receptors, this clustering is usually driven by binding to antibodies of differing affinity in turn bound to multivalent antigen. As a consequence of this activation mechanism, accounting for and rationally manipulating IgG effector function is complicated by, among other factors, the contribution of differing affinities to multiple \fcgr{}s and changes in the valency of antigen binding. In this study, we show that a model of multivalent receptor-ligand binding can effectively account for the contribution of IgG-\fcgr{} affinity and immune complex valency. This model in turn enables us to make specific predictions about the effect of immune complexes of defined composition. In total, these results enable rationally designed IgG effector function, or deconvolution of function, in both a forward and reverse manner.
\end{abstract}
\begin{document}

\flushbottom
\maketitle

\thispagestyle{empty}

\section{Introduction}

% This paragraph addresses the idea that IgGs are relevant to disease
Antibodies are critical and central regulators of the immune response. Antibodies of the IgG isotype interact with FcgR receptors expressed widely on innate immune effector cells. Regulation of effector cell function is a critical component of the IgG therapy's use in cancer and autoimmune diseases. Effector function operates through multiple cell types---including macrophages, monocytes, neutrophils, and NK cells---and multiple process such as antibody-dependent cell-mediated cytotoxicity (ADCC), promoting antigen presentation, and cytokine response. In addition to their effect in isolation, IgG therapies can show synergistic effect in cancers in combination with checkpoint and cytokine-mediated immunotherapies \ac{Moynihan:2016jb, Zhu:2015gy}. These immunotherapeutic effects in combination with antibodies' ability to operate as signaling modulators through competitive binding and opsonization make IgG biologic agents particularly versatile therapeutic agents.

% This paragraph addresses the idea that efforts have been made to manipulate IgG function
An ability to quantitatively predict IgG-\fcgr{} function would aid understanding and treatment of cancer and autoimmune diseases. Efforts to engineer improved effector responses to IgG treatment include mutated Fc variants with biased \fcgr{} binding, deglycosylated Fc domains with the effect of modulating \fcgr{} binding, and alternative IgG subclasses with distinct binding profiles. In addition, particularly in cases where antigen and antibody are exogenously provided, avidity and binding affinities may be manipulated coordinately. % These have generally had X limitations through...

% This paragraph addresses the idea that people have tried to understand IgG function
Multiple efforts have sought to improve our understanding of IgG-mediated effector function. These include efforts to carefully quantify the individual, monovalent \fcgr{}-IgG affinities \ac{Bruhns:2009kg}. In addition, previous studies have characterized the effects of IgG glycosylation (which serves to modulate \fcgr{} affinity) and immune complex avidity on the binding of IgG-antigen complexes \ac{Lux:2013iv}. Genetic models have made it possible to remove certain \fcgr{}s and examine the consequent effect on IgG treatment in various cancer models\cite{Clynes:2000ga}. Finally, comparison of antibodies with matched variable regions but differing Fc domains has allowed the effects to be compared, through with necessarily pleiotropic differences between each of these classes \ac{Nimmerjahn:2005hu}.

% This paragraph addresses the idea that quantitative models have been used to understand immune receptors
Models of multivalent ligand binding to monovalent receptors have been successfully employed to study function of other immune receptors with known, corresponding binding models\ac{Perelson:1980fs, Perelson:1980ds, Hlavacek:1999gb}. For example, a two-component binding model can capture the effect of T cell receptor activation \ac{Stone:2001fm}. However, the \fcgr{} family presents the considerable additional challenge of multiple distinct receptor classes being expressed simultaneously within cells. Additionally, the multiple \fcgr{}s present, with activating and inhibitory roles, ensure that any manipulation of immune complex composition will necessarily have multivariate effects. The same challenge exists for other paired receptor-ligand families, including other immunoglobulin classes and the many phosphatidylserine receptors. Thus, while the underlying theoretical models of multivalent binding are long-standing, \fcgr{}-IgG interactions are especially suited for developments in the ability to rigorously link these models to experimental observations.

In this study, we have employed a model of multivalent immune complex binding to FcgR receptors and show that it can capture the experimentally measured binding at differing valencies. Applying this model, we show it can quantitatively predict effector response in response to diverse interventions in a forward manner, and can deconvolve the causal factors of response in a reverse fashion. More broadly, these results demonstrate the power of a unified binding model tied to computational inference techniques linking theory and experimental observation.


%\section{Methods}

\subsection{Data}

All binding measurements were generously provided by Anja Lux of the Nimmerjahn Lab. Binding was measured using ELISA. The quantity of ligand bound in the original data is measured in terms of mean fluorescent intensity (MFI). Four replicates were used in the lab's assays. We normalized the data by replicate in our simulations.

% TODO: The binding data will need to be described in the level of detail Lux et al has since we will be including original data.

\subsection*{Equations}

The equilibrium binding of TNP-X-BSA to Fc\textgamma{}Rs was modeled using an equilibrium model of multivalent ligand binding to monovalent receptors expressed uniformly on a cell surface presented by Stone \textit{et al}. (2001)\textsuperscript{b}. The model assumes that all cells are submerged in a solution with a saturating concentration of ligand. 
% TODO: Does this really assume a saturating concentration of ligand?
According to the model, the number of ligand bound $i$-valently to the cell at equilibrium, $v_{eq}$, can be found using the relation
\begin{equation}
v_{i,eq} = {v\choose i} (K_X)^{i-1} \frac{L_0}{K_D} (R_{eq})^i.
\end{equation}
Here, $K_X$ is a cross-linking constant with units of \# per cell\textsuperscript{b} (for our purposes, $K_X$ is assumed to be identical for all combinations of Fc\textgamma{}R and IgG,) $L_0$ is the concentration of ligand, and $R_{eq}$ is the number of receptors bound at equilibrium.\textsuperscript{c} Consequently,
\begin{equation}
L_{bound} = \sum_{i=1}^{v} v_{i,eq} = \sum_{i=1}^{v} {v\choose i} (K_X)^{i-1} \frac{L_0}{K_D} (R_{eq})^i,
\end{equation}
where $L_{bound}$ is the total number of ligand bound at equilibrium and $v$ is the effective avidity of the ligand. It is important to distinguish between the effective avidity and the actual avidity of a multivalent ligand, as steric effects might prevent such a ligand from binding to $v_{actual}$ receptors at once. We will later show that this was the case in Lux's experiments. $R_{eq}$ is a function of $v$, $L_0$, $K_D$, $K_X$, and $Rtot$, the total number of receptors expressed on the cell surface, and can be approximated numerically using the following relationship when these parameters are known:
\begin{equation}
R_{tot} = R_{eq} \Big(1+v \frac{L_0}{K_D} (1+K_X R_{eq})^{v-1}\Big).
\end{equation}
Let $R_{eq}(R_{tot},v,L_0,K_D,K_X)$ be the numerical approximation of $R_{eq}$ given the parameters listed. Consequent of Equation (1), $R_{multi}$, the number of receptors that are clustered with at least one other receptor at equilibrium, can be found as follows:
\begin{equation}
R_{multi} = \sum_{i=2}^{v} iv_{i,eq}.
\end{equation}


\subsection{Calculations}

\subsubsection{Parameters and Assumptions}

Association constants for all combinations of IgG and Fc\textgamma{}R were obtained from Bruhns et al. (2009), and corresponding dissociation constants ($K_D$) were calculated using these.\textsuperscript{a} In each replicate of Lux's assay, cells were coincubated with 5 \textmu{}g/ml TNP-X-BSA. Because the molar masses of a 2,4,6-trinitrophenyl groups and of BSA are approximately 173 Da and 66430 Da, respectively, we represented the molar concentrations of TNP-4-BSA and TNP-26-BSA as 74~nM and 70~nM.\textsuperscript{f} For the sake of parsimony, we assumed that the cross-linking constant $K_X$ is identical for all combinations of Fc\textgamma{}R and IgG. We also assumed that there was a conversion factor between the number of ICs bound and MFI for both TNP-4-BSA and TNP-26-BSA and that these two conversion factors were likely different. Lastly, we assumed that, due to steric effects, the effective avidities of TNP-4-BSA and (especially) TNP-26-BSA might be different that their actual avidities. This required us, in total, to fit 11 parameters: the total expression level $R_{tot}$ for each Fc\textgamma{}R, $K_X$, conversion factors from ligand bound to MFI for both TNP-BSAs ($c_{4}$ and $c_{26}$, respectively), and effective avidities for both TNP-BSAs ($f_{4}$ and $f_{26}$, respectively).

Based on the order of magnitude of cross-linking constant values presented in the aforementioned paper by Stone et al., we assumed that our cross-linking constant $K_X$ should be between 

\subsubsection*{Model Fitting}

We fit the aforementioned parameters by minimizing a log-likelihood function of a form similar to that presented by Kitagawa (1986).\textsuperscript{g} We assume that the standard error of MFI for any combination of Fc\textgamma{}R, IgG, and IC (TNP-4-BSA or TNP-26-BSA) is proportional to the mean of the normalized samples from trials involving that combination of conditions. Therefore, as opposed to fitting a standard error $\sigma$ for each combination of Fc\textgamma{}R, IgG, and IC, we fit a standard error coefficient $\sigma^*$ for each combination, which, when multiplied by the mean of the corresponding trial sample, equals the standard deviation of the error for that combination. Therefore, we assume that the likelihood function of any set of parameters granted a specific combination $i$ of Fc\textgamma{}R, IgG, and IC is of the form:

\begin{equation}
\mathcal{L}(K_X,c_4,c_{26},f_4,f_{26},\sigma^*|i) = \prod_{j=1}^{n_i}{(\sigma^*\mu_i\sqrt{2\pi})^{-1}\exp{\Big(\frac{(y_{ij}-c_iL)^2}{2(\sigma^*\mu_i)^2}\Big)}},
\end{equation}

where $n_i$ is the number of samples from the assay for that particular $i$, $y_{ij}$ is the $j$\textsuperscript{th} normalized MFI from the sample representing $i$, and $\mu_i$ is the algebraic mean of the sample representing $i$. The total log-likelihood, then, for particular a set of parameters is:

\begin{equation}
\begin{split}
l_{total}=\sum_i\log(\mathcal{L}(R_{tot,Fc\gamma RIIIA-158V},K_X,c_4,c_{26},f_4,f_{26},\sigma^*|i))
\end{split}
\end{equation}
 
We integrated under the log-likelihood distribution $l$ over all twelve parameters listed above using the mhsample implementation of the Metropolis-Hastings algorithm in MATLAB.\textsuperscript{e,h} We sampled 10 million points with using a burn-in length of 1000.

\section{Results}


\subsection{IgG-FcgR binding varies with affinity and avidity}



\begin{figure}[t]
  \centering
  \includegraphics[width=5in]{./Figures/Fig1-BindingData.pdf}
  \caption{\fcgr{} binding changes with FcgR-IgG pair and avidity. }
  \label{Binding}
\end{figure}




\subsection{A multivalent binding model accounts for variation in IgG-FcgR binding}





\begin{figure}[t]
  \centering
  \includegraphics[width=5in]{./Figures/Fig1-BindingData.pdf}
  \caption{Regression of the multivalent binding model to the \fcgr{} binding data.}
  \label{ModelFit}
\end{figure}






\subsection{}











\subsection{}






%\section{Discussion}




%\section{Acknowledgements}

This work was supported by NIH DP5-OD019815 to A.S.M. \textbf{Competing financial interests:} The authors declare no competing financial interests.

\section{Author contributions statement}

A.S.M. conceived the experiment(s),  A.S.M. and R.A.R. conducted the experiment(s), A.S.M. and R.A.R. analyzed the results.  All authors reviewed the manuscript.

\printbibliography

%a -> Lux 2013
%b -> Stone 2001
%c -> Perelson 1981
%d -> Bruhns
%e -> Paper presenting Metropolis-Hastings algorithm for the first time
%%f -> https://www.biosearchtech.com/tnpbsa-bovine-serum-albumin
%g -> \textit{Akaike Information Criterion Statistics}
%h -> mhsample documentation from MATLAB?

%\cleardoublepage
\section{Supplement}


\begin{table}
\begin{center}
\begin{tabular}{|c|c|c|c|c|c|c|c|c|c|}
\hline
Fc\textgamma{}R & IgG & \multicolumn{8}{c|}{Replicate} \\
\cline{3-10}
& & \multicolumn{4}{c|}{TNP-4-BSA} & \multicolumn{4}{c|}{TNP-26-BSA} \\
\cline{3-10}
& & 1 & 2 & 3 & 4 & 1 & 2 & 3 & 4 \\
\hline
\multirow{5}{5em}{Fc\textgamma{}RIA} & background & 113 & 5 & 7 & 8 & 38 & 6 & 8 & 10 \\
& IgG1 & 441 & 427 & 784 & 172 & 318 & 552 & 900 & 248 \\
& IgG2 & 440 & 578 & 997 & 176 & 404 & 392 & 894 & 124 \\
& IgG3 & 442 & 407 & 905 & 185 & 322 & 515 & 957 & 242 \\
& IgG4 & 388 & 428 & 618 & 139 & 530 & 431 & 810 & 172 \\
\hline
\multirow{5}{5em}{Fc\textgamma{}RIIA-131R} & background & 14 & 8 & 9 & 9 & 18 & 9 & 12 & 12 \\
& IgG1 & 130 & 177 & 274 & 58 & 266 & 604 & 1016 & 281 \\
& IgG2 & 50 & 23 & 57 & 13 & 149 & 139 & 381 & 48 \\
& IgG3 & 184 & 193 & 408 & 87 & 523 & 621 & 1297 & 281 \\
& IgG4 & 44 & 21 & 29 & 17 & 265 & 151 & 343 & 72 \\
\hline
\multirow{5}{5em}{Fc\textgamma{}RIIA-131H} & background & 38 & 4 & 4 & 5 & 53 & 5 & 5 & 10 \\
& IgG1 & 327 & 308 & 539 & 113 & 660 & 585 & 1298 & 377 \\
& IgG2 & 277 & 203 & 174 & 36 & 533 & 394 & 742 & 156 \\
& IgG3 & 307 & 334 & 305 & 135 & 421 & 674 & 1173 & 405 \\
& IgG4 & 148 & 54 & 6 & 20 & 630 & 517 & 1005 & 145 \\
\hline
\multirow{5}{5em}{Fc\textgamma{}RIIB} & background & 7 & 6 & 7 & 4 & 9 & 7 & 9 & 10 \\
& IgG1 & 80 & 200 & 262 & 30 & 139 & 737 & 1108 & 230 \\
& IgG2 & 17 & 27 & 40 & 8 & 63 & 111 & 260 & 32 \\
& IgG3 & 124 & 231 & 349 & 62 & 442 & 942 & 746 & 244 \\
& IgG4 & 33 & 44 & 100 & 21 & 183 & 331 & 605 & 68 \\
\hline
\multirow{5}{5em}{Fc\textgamma{}RIIIA-158F} & background & 7 & 6 & 8 & 6 & 7 & 7 & 8 & 11 \\
& IgG1 & 262 & 676 & 1015 & 269 & 709 & 1333 & 1935 & 698 \\
& IgG2 & 14 & 31 & 31 & 14 & 23 & 26 & 68 & 19 \\
& IgG3 & 328 & 528 & 790 & 267 & 788 & 1668 & 2628 & 891 \\
& IgG4 & 9 & 8 & 16 & 10 & 17 & 21 & 46 & 20 \\
\hline
\multirow{5}{5em}{Fc\textgamma{}RIIIA-158V} & background & 49 & 8 & 14 & 7 & 49 & 8 & 9 & 8 \\
& IgG1 & 295 & 596 & 1016 & 334 & 488 & 813 & 1330 & 383 \\
& IgG2 & 105 & 186 & 293 & 65 & 177 & 362 & 573 & 114 \\
& IgG3 & N/A & 623 & 1055 & 365 & 924 & 1317 & 1861 & 676 \\
& IgG4 & 122 & 136 & 262 & 91 & 313 & 558 & 1026 & 239 \\
\hline

\end{tabular}
\caption{Raw binding measurements. The original MFIs from the Nimmerjahn Lab's experiment, normalized by replicate. No data was recorded in Replicate 2 for the combination of Fc\textgamma{}RIIIA-158V, IgG3, and TNP-4-BSA. }
\end{center}
\end{table}

\end{document}